\documentclass{myclass}
\usepackage{verbatim}
\usepackage{array}
\usepackage{listings}
\usepackage{fancyvrb}

\usepackage[utf8]{inputenc}
\usepackage[T1]{fontenc}
\usepackage{textcomp}
\usepackage{multicol}
\usepackage{mathtools}
\usepackage{amsmath}
\usepackage{wrapfig}
\usepackage{amssymb}
\usepackage{amsmath,amsfonts,amssymb,amsthm,epsfig,epstopdf,titling,url,array}
\usepackage{hyperref}
\usepackage{eso-pic}
\usepackage{pgf}
\usepackage{tikz}
\usepackage{graphicx}

% figure support
\usepackage{import}
\usepackage{xifthen}
\pdfminorversion=7
\usepackage{pdfpages}
\usepackage{transparent}
\usepackage{xcolor}

%formatting
\setlength{\parindent}{0em}
\setlength{\parskip}{1em}
%\renewcommand{\baselinestretch}{2.0}

\newtheorem*{definition}{Definition}
\newtheorem*{theorem}{Theorem}

\newcommand{\incfig}[1]{%
  \center
    \def\svgwidth{0.9\columnwidth}
    \import{./figures/}{#1.pdf_tex}
}


\newcommand{\incimg}[1]{%
\center
\includegraphics[width=0.9\columnwidth]{images/#1}
}


\pdfsuppresswarningpagegroup=1

\title{a}

\begin{document}
\maketitle
\tableofcontents
\newpage

\section{First section}
This is the first paragraph, contains some text to test the paragraph
interlining, paragraph indentation and some other features. Also, is 
easy to see how new paragraphs are defined by simply entering a double 
blank space.

Third phrase

\begin{minipage}[t]{0.7\textwidth}
lañkjdflñkdsjfñlkajf alkdjfladf lañdkjf aldkjf aldfj adlfkjladf ladfj dk adka akdj ad adf ladfj ladf   ldfj j aldj.

Hola

\end{minipage}
\begin{minipage}{0.3\textwidth}
\incfig{test1}
\end{minipage}

\begin{minipage}[t]{0.3\textwidth}
\incfig{test1}
\end{minipage}
\begin{minipage}[t]{0.3\textwidth}
Hola
\end{minipage}


\begin{definition}
first definition
\end{definition}

\begin{proof}
$0 
\end{proof}

\end{document}
